% STATUS: DRAFT
\chapter{Open problems across all parts}\label{ch:open_all}

\begin{physmotiv}
A map of what is not known, organized by the Parts of this book.
\end{physmotiv}

\begin{description}[leftmargin=1em]
\item[Parts II--III (algebraic foundations).] Classification of factors beyond the hyperfinite case is wide open; which type III$_1$ factors arise as local algebras of interacting QFTs in $d=4$? Does the hyperfinite property (injectivity) of local algebras hold in general, or only under nuclearity assumptions?
\item[Part IV (modular theory).] Explicit modular Hamiltonians for generic regions and states in interacting QFT are known only in special cases (wedges, balls in CFT, a few free-field examples). Is there a geometric characterization of when the modular flow is local?
\item[Part VI (AQFT).] Existence of interacting nets in $3+1$ dimensions; the relation between algebraic and constructive QFT.
\item[Part VII (de Sitter and observers).] See \cref{ch:open_ds}: non-perturbative definition, observer dependence, the additive constant, multi-observer algebras, dynamics.
\item[Part VIII (black holes).] The finite-$N$ origin of the large-$N$ type III$_1$ limit; the algebraic description of black hole evaporation including the island formula; the type of the algebra of the exterior of an evaporating black hole.
\item[Part IX (noncommutative geometry).] A Lorentzian spectral triple formalism; the status of the spectral Standard Model; the relationship, if any, to the gravitational algebras of Part VII.
\end{description}

\begin{keyidea}
The unifying open question: \emph{what is the algebra of observables of quantum gravity, and which of its properties (type, trace, entropy) are robust beyond the semiclassical limit?}
\end{keyidea}
